% HD-F1-T32. GPT-6 Astra Pro, /root/transport; 22 November 2026 UTC.
\section[HD-F1-T32: effective linear and unitary level reduction]{HD-F1-T32: an effective natural-dimension cutoff leaves only level one over the fields of size two and three}
\label{sec:f1-large-dimension-character-level-reduction}

Every remaining T2-OTHER source whose almost simple projective
coefficient image has socle \(PSL_N(q)\) or \(PSU_N(q)\) satisfies
the following explicit alternative:
\[
 \boxed{\quad N<10^{62}\quad\text{or}\quad
                  (q,j)\in\{(2,1),(3,1)\}.\quad}
 \tag{F1T32.1}
\]
The defining field already satisfies \(q\le27\) by T18--T19.
For the second alternative, level means the level of the
actual constituent inflated along the universal central-cover
map below. All these constituents have the same level, even
when the restriction is reducible and original inertia induces
outer automorphisms. This is an effective reduction of the
original source gate; it does not assert the existence of a
source in either retained alternative.

\paragraph{1. The actual perfect layer and its lifted automorphisms.}
Retain the original faithful absolutely irreducible
\(G\subset GL_K(V)\), \(m=\dim_KV\), its scalar centre
\(Z\), and its projective socle \(S\). Let \(P\) be the
inverse image of \(S\) and let \(L=P'\). The actual maps
of T16 give
\[
 P=ZL,\qquad L=[L,L],\qquad Z(L)=L\cap Z,
 \qquad V_{\mathbf C}^{L}=0.
 \tag{F1T32.2}
\]
It remains to analyze \(N\ge N_0:=10^{62}\).
Put \(H=SL_N(q)\) or \(SU_N(q)\), respectively.
The perfectness and universal-cover exceptions of
Malle--Testerman, Theorem~24.17 and
Remark~24.19/Table~24.3, have \(N\le6\) in these two
families. Thus \(\pi:H\to S\) is the universal central
cover. There is a unique comparison homomorphism
\[
 \beta:H\longrightarrow L,\qquad
       (L\longrightarrow S)\beta=\pi.
 \tag{F1T32.3}
\]
Its image \(I\) satisfies \(L=Z(L)I\), so perfectness
gives \(L=[I,I]\subset I\); hence \(\beta\) is
surjective, with central kernel.

For \(g\in G\), denote its conjugation action on \(L\)
by \(c_g\) and its induced automorphism of \(S\) by
\(a_g\). Apply the universal property to the central
extension \(a_g^{-1}\pi:H\to S\). It gives a unique map
\(\widetilde a_g:H\to H\) satisfying
\(\pi\widetilde a_g=a_g\pi\). The corresponding map
for \(a_g^{-1}\), together with uniqueness over the
identity, proves that \(\widetilde a_g\) is an
automorphism. Applying the same uniqueness to the central
extension \(a_g^{-1}(L\to S)\) gives the exact square
\[
                  \beta\widetilde a_g=c_g\beta.
 \tag{F1T32.4}
\]
Thus the automorphism on every inflated constituent is
the one induced by the original matrix \(g\).

Write \(\epsilon=1\) in the linear case and \(-1\) in
the unitary case, and put \(Q=q\) and \(q^2\), respectively.
On \(H\), restrict the defining reducible Weil character
\[
 \tau(h)=\epsilon^N(\epsilon q)^{d(h)},
 \qquad d(h)=\dim_{\mathbf F_Q}\ker(h-1).
 \tag{F1T32.5}
\]
Guralnick--Larsen--Tiep, \emph{Forum of Mathematics, Pi}
\textbf{8} (2020), e2, Definition~2 and Lemma~7.1,
identify the level of \(\varphi\in\operatorname{Irr}(H)\)
with the least integer \(j\ge0\) for which
\(\varphi\) occurs in \(\tau^j\).

The exact automorphism inputs are Malle--Testerman
Theorem~24.24 for the split case and Steinberg,
\emph{Lectures on Chevalley Groups}, Theorem~36 for the
twisted case. Their generators lift on the displayed
natural module to conjugation by an invertible matrix,
entrywise field automorphisms, and, in the split case,
inverse transpose. These preserve \(d(h)\): conjugation
identifies the kernels; a field automorphism gives the
corresponding semilinear bijection of kernels; and
\[
 \operatorname{rank}(h^{-t}-1)
 =\operatorname{rank}\bigl(-h^{-t}(h^t-1)\bigr)
 =\operatorname{rank}(h-1).
 \tag{F1T32.6}
\]
The uniqueness just proved identifies these standard
lifts with the lifts in (F1T32.4). Consequently every
\(\widetilde a_g\) fixes \(\tau\), and hence fixes
the level of each constituent.

For completeness, the isotypic summands of
\(V_{\mathbf C}|_L\) are permuted by \(G\).
The sum over any one orbit is \(G\)-stable, so absolute
irreducibility makes this action transitive.
Equation (F1T32.4) and the preceding level invariance
therefore give one common level \(j\) on the actual
inflated restriction
\[
 V_{\mathbf C}|_H=\bigoplus_i a_i W_i,
 \qquad a_i\in\mathbf Z_{>0},\qquad
 m=\sum_i a_i\dim W_i.
 \tag{F1T32.7}
\]
Every \(W_i\) is nontrivial by (F1T32.2), so \(j\ge1\).

\paragraph{2. Precise degree and uniform-support inputs.}
The published Guralnick--Larsen--Tiep Theorem~1.3,
restated and proved as Theorem~7.3, gives, for an
irreducible \(H\)-character \(\varphi\) of level \(j\),
\[
 \varphi(1)\ge\frac{q^{j(N-j)}}{2(q+1)},
 \qquad
 j\ge N/2\ \Longrightarrow\
 \varphi(1)>\frac23q^{N^2/4-3}.
 \tag{F1T32.8}
\]
The first bound uses the weaker of its two denominators
and so applies to both groups. In this source gate
\(2(q+1)\le56\).

Larsen--Tiep, \emph{Annals of Mathematics}
\textbf{200} (2024), 1--70, Theorem~5.5 in the inspected
author version \texttt{arXiv:2403.09046v1}, proves
\[
 \frac{|\varphi(h)|}{\varphi(1)}
 \le\varphi(1)^{-\sigma\operatorname{supp}(h)/N},
 \qquad
 \sigma=\frac1{6507\cdot2^{21}\cdot10^{18}}>10^{-29}.
 \tag{F1T32.9}
\]
Here support is the codimension of a largest eigenspace
on the natural module over an algebraic closure. Thus
every noncentral \(h\in H\) has support at least one.
The exact theorem covers both \(SL_N(q)\) and
\(SU_N(q)\), all irreducible complex characters and
all elements. We use the stated effective constant,
not an unspecified asymptotic constant.

Set
\[
             J=10^{30},\qquad N_0=100J^2=10^{62},
             \qquad \sigma_0=10^{-29}.
 \tag{F1T32.10}
\]
For \(J<j\le N/2\), (F1T32.8) gives
\[
 \frac{\log_2\varphi(1)}N
 \ge j\left(1-\frac jN\right)-\frac6N
 \ge\frac j2-\frac6N.
 \tag{F1T32.11}
\]
Its product with \(\sigma_0\) exceeds four, since
\(j>J\) and \(N\ge100J^2\).
For \(j\ge N/2\), the second part of (F1T32.8)
gives instead
\[
 \frac{\log_2\varphi(1)}N>\frac N4-\frac4N,
 \tag{F1T32.12}
\]
whose product with \(\sigma_0\) also exceeds four.
Combining these estimates with (F1T32.9) gives the
fully effective large-level bound
\[
       \frac{|\varphi(h)|}{\varphi(1)}<\frac1{16}
       \qquad(j>J,\ h\notin Z(H)).
 \tag{F1T32.13}
\]

\paragraph{3. The actual dual-pair projector improves the low-level constants.}
Let \(2\le j\le J\). Write
\(\mathcal G=GL_N(q)\) or \(GU_N(q)\), and
\(\mathcal B=GL_j(q)\) or \(GU_j(q)\), respectively.
By Lemma~7.1 and Corollary~8.6 of the published
Guralnick--Larsen--Tiep paper, \(\varphi\) has a
unique extension \(\chi\in\operatorname{Irr}(\mathcal G)\)
of true level \(j\), and \(\chi|_H=\varphi\).
The ranges used in its Section~8 hold because
\[
 \frac32(j^2+j+1)\le N-1,
 \qquad j^2+j\le N-1.
 \tag{F1T32.14}
\]

The relevant representation is the reducible Weil
representation on the natural tensor product
\(\mathbf F_Q^N\otimes\mathbf F_Q^j\).
For \(\alpha\in\operatorname{Irr}(\mathcal B)\),
the actual multiplicity-space projector has character
\[
 D_\alpha(h)=\frac1{|\mathcal B|}
       \sum_{s\in\mathcal B}
           \tau_{Nj}(h\otimes s)\overline{\alpha(s)}.
 \tag{F1T32.15}
\]
Equivalently it is obtained from
\(\operatorname{Hom}_{\mathcal B}
 (W_\alpha,\mathcal W)\), where \(\mathcal W\)
is that Weil representation; the standard central
idempotent first selects its \(\alpha\)-isotypic
summand. The source proves that
\(D_\alpha=D_\alpha^\circ+D_\alpha^0\), where
\(D_\alpha^\circ=\chi\) and \(D_\alpha^0\)
is an actual character of lower true level.
For the present ranges this is Theorem~1.1 and the
complete arguments in Sections~8.1 and 8.3.
No substitution of this tensor representation for the
original coefficient local system is made.

The exact inequalities (8.11) and (8.18) of that proof,
obtained from (F1T32.15), yield
\[
 \frac{|\chi(h)|}{\chi(1)}
 \le\frac{q^{-j}+E_\epsilon}{1-E_\epsilon}
       \qquad(h\in\mathcal G\setminus Z(\mathcal G)),
 \tag{F1T32.16}
\]
where the following upper bounds for the common
numerator error and denominator error suffice:
\[
 \begin{aligned}
 E_+&=|GL_j(q)|q^{-N}
        \left(1+\sqrt{8q^{j^2+j}}\right),\\
 E_-&=|GU_j(q)|q^{-N}
        \left(1+\sqrt{\tfrac{413}{25}q^{j^2+j-1-2N}}\right).
 \end{aligned}
 \tag{F1T32.17}
\]
Indeed the exceptional term in the numerator of each
source inequality is at most \(q^{(N-1)j}\); its
remaining non-remainder term is at most
\(|\mathcal B|q^{N(j-1)}\), since
\((N-1)(j-1)+1\le N(j-1)\).
Dividing numerator and denominator by \(q^{Nj}\)
gives precisely (F1T32.16)--(F1T32.17).
The character remainder is estimated in source
equations (8.10) and (8.17), respectively.

The elementary group-order bounds in that proof are
\(|GL_j(q)|<q^{j^2}\) and
\(|GU_j(q)|\le\tfrac32q^{j^2}\).
Since \(\sqrt8<3\), \(\sqrt{413/25}<5\), and
\(j\ge2\), both displayed errors satisfy
\[
 \begin{aligned}
 E_\epsilon
 &\le9q^{(3j^2+j)/2-N}
 \le9q^{2J^2-N}
 \le\frac9{2^{N-2J^2}}
 <\frac1{100}.
 \end{aligned}
 \tag{F1T32.18}
\]
For the last inequality, \(N-2J^2\ge98J^2\ge10\)
and \(900<2^{10}\). Consequently
\[
 \frac{|\varphi(h)|}{\varphi(1)}
 \le\frac{1/4+1/100}{1-1/100}
 =\frac{26}{99}<\frac3{10}
 \qquad(2\le j\le J,\ h\notin Z(H)).
 \tag{F1T32.19}
\]
Thus the integral effective cutoff excludes level two
over \(q=2\) as well as every larger level; use of
the coarser headline constants alone would miss it.

\paragraph{4. Level one and its remaining field sizes.}
For \(j=1\), the published Section~8 proofs give in
both cases the exact estimates
\[
 \chi(1)\ge\frac{q^N-q}{q-\epsilon},\qquad
 |\chi(h)|\le\frac{q^{N-1}+q}{q-\epsilon}
       \quad(h\notin Z(\mathcal G)).
 \tag{F1T32.20}
\]
The same irreducible extension and restriction map
applies. For \(q\ge4\), put \(t=q^{1-N}\).
Since \(N\ge N_0\), \(t<1/100\), and hence
\[
 \frac{|\varphi(h)|}{\varphi(1)}
 \le\frac{q^{-1}+t}{1-t}
 \le\frac{26}{99}<\frac3{10}.
 \tag{F1T32.21}
\]
Together with (F1T32.13) and (F1T32.19), this bounds
every noncentral element by the uniform ratio \(3/10\)
unless the common level is one and \(q\in\{2,3\}\).

\paragraph{5. Return the bounds to the original inertia maps.}
Take any \(h\in L\setminus Z(L)\). Its preimages
under \(\beta\) are noncentral in \(H\), because
their images in \(S\) are nontrivial.
Outside the retained level-one cases, apply the bounds
just proved to every actual summand in (F1T32.7) and
add with the actual multiplicities. This gives
\[
 |\operatorname{Tr}_V(h)|\le\frac3{10}m,
 \qquad
 \operatorname{rank}(h-1)\ge\frac7{20}m.
 \tag{F1T32.22}
\]
The second inequality follows because finite order
makes all eigenvalues unit complex numbers:
\(m-2\operatorname{rank}(h-1)\le
 |\operatorname{Tr}_V(h)|\).

For every nonscalar original \(g\in G\), choose
\(e\in L\) with \([g,e]\notin Z(L)\) by the actual
projective centralizer argument of T16. The matrix map
\[
 [g,e]-1=
   \bigl((g-\lambda)e-e(g-\lambda)\bigr)g^{-1}e^{-1}
 \tag{F1T32.23}
\]
has rank at most \(2\operatorname{rank}(g-\lambda)\)
for every complex scalar \(\lambda\). Therefore every
nonscalar original inertia has moving rank at least
\(7m/40\). A nonidentity scalar inertia has moving
rank exactly \(m\), so obeys the same bound.
The unchanged source identities are
\[
 2m+6=\sum_{\nu=1}^b\operatorname{rank}(g_\nu-1),
 \qquad b\ge12.
 \tag{F1T32.24}
\]
They imply
\[
          2m+6\ge\frac{21}{10}m,
          \qquad m\le60.
 \tag{F1T32.25}
\]
Every summand in (F1T32.7), however, has degree
greater than sixty. For \(1\le j\le N/2\),
\(j(N-j)-(N-1)=(j-1)(N-j-1)\ge0\), so (F1T32.8) gives
\(\dim W_i\ge2^{N-1}/56>60\).
For \(j\ge N/2\), its second bound gives the same
strict conclusion. Already \(N\ge16\) suffices for
both inequalities. Since \(m\ge\dim W_i\), this
contradicts (F1T32.25), proving (F1T32.1).

\paragraph{6. Original geometric and realization maps; source scope.}
The proof uses the actual coefficient inflation
(F1T32.3), its commuting automorphism square
(F1T32.4), the multiplicity-space projector
(F1T32.15) for numerical character estimates, and the
original matrix identity (F1T32.23). The map of the
original cover \(C\to\mathbf P^1\) is its identity;
the middle-extension coefficient sheaf, branch stalks
\(V^{g_\nu}\), parabolic cohomology and its cup form
have their identity maps. The rational primitive
idempotent, its integer-cleared actual Jacobian image,
the prescribed \(K\)-corner and its restricted theta
form are retained. Hence the induced maps of Poincare
biextensions, Hodge filtrations, period differentials
and the Grothendieck/Adams realizations of F01--F04
are identities, with polarization multiplier one.
No universal-cover lift of the curve is asserted.
No new detector or determinant component is counted.

The published GLT source was inspected in its Cambridge
PDF, including the complete degree proof of Section~5,
Lemma~7.1, Theorem~7.3, the full low-level proofs of
Section~8, and Corollary~8.6. The Larsen--Tiep proof
was inspected in the authors' 14 March 2024 version,
including the complete printed proofs of Theorem~4.4,
Propositions~5.2--5.4 and Theorem~5.5; the Annals publication
record was separately verified. The source audit
distinguishes that author PDF from a publisher PDF.
The automorphism classification is the named primary
input: the complete supplied split proof in
Malle--Testerman and the complete supplied twisted
proof sketch in Steinberg were read. The exact lift,
level-invariance, finite cutoff and original-source
rank deductions are proved above.
